% Part: first-order-logic
% Chapter: syntax-and-semantics
% Section: unique-readability

\documentclass[../../../include/open-logic-section]{subfiles}

\begin{document}

\olfileid{fol}{syn}{unq}

\olsection{একক পাঠযোগ্যতা}

\begin{explain}
আমরা যেভাবে !!{formula}s সংজ্ঞায়িত করেছি, তাতে প্রত্যেকটি !!{formula}-এর
\emph{একক পাঠ} নিশ্চিত হয়। অর্থাৎ !!{formula}s-এর গঠনের নিয়ম অনুযায়ী
তাকে নির্মাণ করার মূলত একটিমাত্র উপায় এবং তাকে “ব্যাখ্যা” করারও একটিমাত্র
উপায় আছে। তা না হলে দ্ব্যর্থক !!{formula}s থাকত—অর্থাৎ একাধিক পাঠ বা
ব্যাখ্যাসম্পন্ন !!{formula}s—যা আমরা স্পষ্টতই এড়াতে চাই। আরও বড় কথা,
এই ধর্মটি না থাকলে পরে দেওয়া আমাদের অধিকাংশ সংজ্ঞা ও প্রমাণই কার্যকর
হবে না।

একক পাঠযোগ্যতা নিশ্চিত করে না—এমন !!{formula}s গঠনের ত্রুটিপূর্ণ নিয়ম
দিলে কী হতো, তা দেখলেই বিষয়টি সম্ভবত সবচেয়ে পরিষ্কার হবে। যেমন, সংযোজকের গঠন-নিয়মে
আমরা বন্ধনী দিতে ভুলে গিয়ে নিচের নিয়মটি অনুমোদন করতে পারতাম:
\begin{quote}
$!A$ ও $!B$ যদি !!{formula}s হয়, তাহলে $!A \lif !B$-ও তাই।
\end{quote}
একটি পরমাণু সূত্র $!D$ থেকে শুরু করে এতে $!D \lif !D$ গঠন করা যেত।
এটি এবং $!D$ নিয়ে তারপর $!D \lif !D \lif !D$ পাওয়া যেত। কিন্তু তা
করার দুটি উপায় আছে:
\begin{enumerate}
\item $!D$-কে $!A$ এবং $!D \lif !D$-কে $!B$ ধরি।
\item $!A$-কে $!D \lif !D$ ধরি এবং $!B$ হলো~$!D$।
\end{enumerate}
সেই অনুযায়ী !!{formula}~$!D \lif !D \lif !D$ “পড়ার”ও দুটি উপায়
থাকে। এটি $!B \lif !C$ রূপের, যেখানে $!B$ হলো $!D$ এবং $!C$ হলো
$!D \lif !D$; কিন্তু এটি \emph{একই সঙ্গে} $!B \lif !C$ রূপেরও,
যেখানে $!B$ হলো $!D \lif !D$ এবং $!C$ হলো~$!D$।

এমন হলে আমাদের সংজ্ঞাগুলি সব সময় কাজ করবে না। যেমন, কোনো সূত্রের
!!{main operator} সংজ্ঞায়িত করার সময় আমরা বলি: $!B \lif !C$ রূপের
সূত্রে $\lif$-এর নির্দেশিত সংঘটনটিই !!{main operator}। কিন্তু উপরের
দুটি ভিন্ন উপায়ে $!D \lif !D \lif !D$ সূত্রটিকে $!B \lif !C$-এর
সঙ্গে মেলানো গেলে এক ক্ষেত্রে $\lif$-এর প্রথম সংঘটনটি, অন্য ক্ষেত্রে
দ্বিতীয় সংঘটনটি !!{main operator} হয়। অথচ আমরা চাই, !!{main operator}
যেন !!{formula}-টির একটি \emph{ফলন} হয়; অর্থাৎ প্রত্যেকটি !!{formula}-এ
!!{main
  operator}-এর ঠিক একটিমাত্র সংঘটন থাকতে হবে।
\end{explain}

\begin{lem}
!!a{formula}~$!A$-তে বাঁ ও ডান বন্ধনীর সংখ্যা সমান।
\end{lem}

\begin{proof}
$!A$ যেভাবে নির্মিত হয়েছে তার উপর আরোহ করে আমরা এটি প্রমাণ করি। এর
জন্য দুটি কাজ দরকার: (ক) প্রথমে প্রমাণ করতে হবে, সব পরমাণু সূত্রেরই
উল্লিখিত ধর্মটি আছে (আরোহের ভিত্তি)। (খ) তারপর প্রমাণ করতে হবে, প্রদত্ত
সূত্রগুলি থেকে নতুন সূত্র নির্মাণের সময় পুরোনো সূত্রগুলির ধর্মটি থাকলে
নতুন সূত্রটিরও তা থাকে।

$l(!A)$ দিয়ে $!A$-তে বাঁ বন্ধনীর সংখ্যা এবং $r(!A)$ দিয়ে ডান বন্ধনীর
সংখ্যা বোঝাই; একইভাবে $l(t)$ ও $r(t)$ দিয়ে পদ~$t$-তে যথাক্রমে বাঁ ও
ডান বন্ধনীর সংখ্যা বোঝাই।

\begin{prob}
    যেকোনো পদ~$t$-এর জন্য $l(t) = r(t)$ প্রমাণ করুন।
\end{prob}

\begin{enumerate}
\tagitem{prvFalse}{\indcase{!A}{\lfalse}{$\indfrm$-এ $0$টি বাঁ এবং $0$টি
    ডান বন্ধনী আছে।}}{}

\tagitem{prvTrue}{\indcase{!A}{\ltrue}{$\indfrm$-এ $0$টি বাঁ এবং $0$টি
    ডান বন্ধনী আছে।}}{}

\item \indcase{!A}{\Atom{R}{t_1,\dots,t_n}}{$l(\indfrm) = 1 + l(t_1) +
  \dots + l(t_n) = 1 + r(t_1) + \dots + r(t_n) = r(\indfrm)$। এখানে
  অনুশীলনী হিসেবে রাখা এই তথ্যটি ব্যবহার করছি যে, যেকোনো পদ~$t$-এর
  জন্য $l(t) = r(t)$।}

\item \indcase{!A}{\eq[t_1][t_2]}{$l(\indfrm) = l(t_1) + l(t_2) =
  r(t_1) + r(t_2) = r(\indfrm)$।}

\tagitem{prvNot}{\indcase{!A}{\lnot !B}{আরোহ-অনুমান থেকে
    $l(!B) = r(!B)$। তাই $l(\indfrm) = l(!B) = r(!B) =
    r(\indfrm)$।}}{}

\item \indcase{!A}{(!B \ast !C)}{আরোহ-অনুমান থেকে $l(!B) =
  r(!B)$ এবং $l(!C) = r(!C)$। তাই $l(\indfrm) = 1 + l(!B) + l(!C) =
  1 + r(!B) + r(!C) = r(\indfrm)$।}

\tagitem{prvAll}{\indcase{!A}{\lforall[x][!B]}{আরোহ-অনুমান থেকে
    $l(!B) = r(!B)$। তাই $l(\indfrm) = l(!B) = r(!B) =
    r(\indfrm)$।}}{}

% Print case for \lexists if prvEx and not prvAll
\tagitem{prvAll,notprvEx}{}{\indcase{!A}{\lexists[x][!B]}{আরোহ-অনুমান
    থেকে $l(!B) = r(!B)$। তাই $l(\indfrm) = l(!B) = r(!B) =
    r(\indfrm)$।}}

% Just say similarly if prvEx and prvAll
\tagitem{notprvAll,notprvEx}{}{\indcase{!A}{\lexists[x][!B]}{একইভাবে।}}
\end{enumerate}
\end{proof}

\begin{defn}[প্রকৃত পূর্বাংশ]
একটি প্রতীকক্রম~$!B$-এর সঙ্গে একটি অশূন্য প্রতীকক্রম সংযুক্ত করলে যদি
প্রতীকক্রম~$!A$ পাওয়া যায়, তাহলে $!B$-কে $!A$-এর একটি
\emph{প্রকৃত পূর্বাংশ} বলে।
\end{defn}

\begin{lem}\ollabel{lem:no-prefix}
$!A$ যদি !!a{formula} এবং $!B$ যদি $!A$-এর প্রকৃত পূর্বাংশ হয়,
তাহলে $!B$ !!a{formula} নয়।
\end{lem}

\begin{proof}
অনুশীলনী।
\end{proof}

\begin{prob}
\olref[fol][syn][unq]{lem:no-prefix} প্রমাণ করুন।
\end{prob}

\begin{prop}
\ollabel{prop:unique-atomic}
$!A$ একটি পরমাণু !!{formula} হলে নিচের শর্তগুলির ঠিক একটি পূরণ করে।
\begin{enumerate}
\tagitem{prvFalse}{$!A \ident \lfalse$।}{}
\tagitem{prvTrue}{$!A \ident \ltrue$।}{}
\item $!A \ident \Atom{R}{t_1,\dots,t_n}$, যেখানে $R$ একটি
  $n$-স্থানীয় !!{predicate}, $t_1$, \dots, $t_n$ পদ এবং $R$,
  $t_1$, \dots, $t_n$-এর প্রত্যেকটি এককভাবে নির্ধারিত।
\item $!A \ident \eq[t_1][t_2]$, যেখানে $t_1$ ও $t_2$ এককভাবে
  নির্ধারিত পদ।
\end{enumerate}
\end{prop}

\begin{proof}
অনুশীলনী।
\end{proof}

\begin{prob}
\olref[fol][syn][unq]{prop:unique-atomic} প্রমাণ করুন। (ইঙ্গিত:
পদের জন্য \olref[fol][syn][unq]{lem:no-prefix}-এর একটি সংস্করণ বিবৃত
ও প্রমাণ করুন।)
\end{prob}

\begin{prop}[একক পাঠযোগ্যতা]
প্রত্যেকটি !!{formula} নিচের শর্তগুলির ঠিক একটি পূরণ করে।
\begin{enumerate}
\item $!A$ পরমাণু।

\tagitem{prvNot}{$!A$ হলো $\lnot !B$ রূপের।}{}

\tagitem{prvAnd}{$!A$ হলো $(!B \land !C)$ রূপের।}{}

\tagitem{prvOr}{$!A$ হলো $(!B \lor !C)$ রূপের।}{}

\tagitem{prvIf}{$!A$ হলো $(!B \lif !C)$ রূপের।}{}

\tagitem{prvIff}{$!A$ হলো $(!B \liff !C)$ রূপের।}{}

\tagitem{prvAll}{$!A$ হলো $\lforall[x][!B]$ রূপের।}{}

\tagitem{prvEx}{$!A$ হলো $\lexists[x][!B]$ রূপের।}{}
\end{enumerate}
তদুপরি, প্রত্যেক ক্ষেত্রে $!B$, অথবা $!B$ ও $!C$, এককভাবে নির্ধারিত।
এর অর্থ, যেমন, এমন আলাদা জোড়া $!B$, $!C$ এবং $!B'$, $!C'$ নেই যাতে
$!A$ একই সঙ্গে
\iftag{prvIf}{$(!B \lif !C)$ ও $(!B' \lif !C')$ রূপের।}{
    \iftag{prvOr}{$(!B \lor !C)$ ও $(!B' \lor !C')$ রূপের।}{
      \iftag{prvAnd}{$(!B \land !C)$ ও $(!B' \land !C')$ রূপের।}{}}}
\end{prop}

\begin{proof}
গঠনের নিয়মগুলি অনুসারে, কোনো !!a{formula} পরমাণু না হলে সেটি অবশ্যই
একটি খোলা বন্ধনী~(, \iftag{prvNot}{$\lnot$,}{} অথবা কোনো পরিমাণসূচক
দিয়ে শুরু হবে। অন্য দিকে, নিচের সংকেতগুলির কোনোটি দিয়ে শুরু হওয়া
প্রত্যেক !!{formula} অবশ্যই পরমাণু: !!a{predicate}, !!a{function},
!!a{constant}\iftag{prvFalse}{, $\lfalse$}{}\iftag{prvTrue}{, $\ltrue$}{}।

তাই আমাদের আসলে শুধু দেখাতে হবে, $!A$ যদি $(!B \ast !C)$ রূপের এবং
একই সঙ্গে $(!B' \mathbin{\ast'} !C')$ রূপের হয়, তাহলে $!B \ident
!B'$, $!C \ident !C'$ এবং $\ast = {\ast'}$।

ধরি $!A \ident (!B \ast !C)$ ও $!A \ident (!B'
\mathbin{\ast'} !C')$ উভয়ই সত্য। তাহলে হয় $!B \ident !B'$, নয়তো তা
নয়। সত্য হলে স্পষ্টতই $\ast = {\ast'}$ এবং $!C \ident !C'$; কারণ
তখন এগুলি $!A$-এর একই স্থান থেকে শুরু হওয়া, সমান দৈর্ঘ্যের
উপপ্রতীকক্রম। অন্য ক্ষেত্রটি হলো $!B \not\ident !B'$। $!B$ ও $!B'$
উভয়ই $!A$-এর একই স্থান থেকে শুরু হওয়া উপপ্রতীকক্রম বলে একটিকে অন্যটির
প্রকৃত পূর্বাংশ হতে হবে। কিন্তু \olref{lem:no-prefix} অনুসারে তা অসম্ভব।
\end{proof}


\end{document}
